\documentclass[11pt]{article}
\usepackage[margin=0.9in]{geometry}
\usepackage{amsmath,amssymb,amsthm,mathtools}
\usepackage[T1]{fontenc}
\usepackage[utf8]{inputenc}
\usepackage{enumitem}
\usepackage{booktabs,longtable,array,seqsplit}
\usepackage{hyperref}
\usepackage{xcolor}
\usepackage{microtype}
\hypersetup{colorlinks=true,linkcolor=blue!45!black,citecolor=blue!45!black,urlcolor=blue!45!black}
\setlength{\emergencystretch}{6em}

\newtheorem{theorem}{Theorem}[section]
\newtheorem{proposition}[theorem]{Proposition}
\newtheorem{lemma}[theorem]{Lemma}
\newtheorem{corollary}[theorem]{Corollary}
\newtheorem{claim}[theorem]{Claim}
\theoremstyle{definition}
\newtheorem{definition}[theorem]{Definition}
\newtheorem{assumption}[theorem]{Assumption}
\theoremstyle{remark}
\newtheorem{remark}[theorem]{Remark}

\newcommand{\R}{\mathbb{R}}
\newcommand{\C}{\mathbb{C}}
\newcommand{\N}{\mathbb{N}}
\newcommand{\Hsp}{\mathcal{H}}
\newcommand{\Fsp}{\mathcal{F}}
\newcommand{\Id}{\mathrm{Id}}
\newcommand{\tr}{\operatorname{tr}}
\newcommand{\diag}{\operatorname{diag}}
\newcommand{\rank}{\operatorname{rank}}
\newcommand{\spec}{\operatorname{spec}}
\newcommand{\spanop}{\operatorname{span}}
\newcommand{\dist}{\operatorname{dist}}
\newcommand{\range}{\operatorname{ran}}
\newcommand{\Ker}{\operatorname{ker}}
\newcommand{\Pinch}{\mathcal{C}}
\newcommand{\Off}{\widetilde{\mathcal{C}}}
\newcommand{\Pperp}{P^{\perp}}
\newcommand{\Qperp}{Q^{\perp}}
\newcommand{\op}{\mathrm{op}}
\newcommand{\F}{\mathrm{F}}
\newcommand{\HS}{\mathrm{HS}}
\newcommand{\ip}[2]{\left\langle #1,#2\right\rangle}
\newcommand{\norm}[1]{\left\lVert #1\right\rVert}
\newcommand{\abs}[1]{\left\lvert #1\right\rvert}
\newcommand{\code}[1]{\nolinkurl{#1}}

\newenvironment{argument}{\paragraph{Argument.}\begin{enumerate}[leftmargin=2em]}{\end{enumerate}}
\newenvironment{proofoutline}{\paragraph{Proof architecture.}\begin{enumerate}[leftmargin=2em]}{\end{enumerate}}
\newcommand{\sourceanchor}[1]{\par\smallskip\noindent\textbf{Source anchor.} #1\par\smallskip}
\newcommand{\sourcecomment}[1]{\begin{remark}#1\end{remark}}
\newenvironment{sourceguide}{\begin{description}[style=nextline,leftmargin=3.3cm,labelwidth=3.1cm]}{\end{description}}
\newenvironment{formalizationmap}{\begin{longtable}{@{}p{0.20\textwidth}p{0.31\textwidth}p{0.40\textwidth}@{}}\toprule Source item & Modern mathematical content & Repository status \\\midrule\endfirsthead\toprule Source item & Modern mathematical content & Repository status \\\midrule\endhead}{\bottomrule\end{longtable}}

\newcommand{\cothh}{\operatorname{coth}}
\newcommand{\sgn}{\operatorname{sgn}}
\newcommand{\Lap}{\mathcal{L}}
\newcommand{\FT}{\mathcal{F}}
\newcommand{\Ioi}{(0,\infty)}
\renewenvironment{formalizationmap}{%
  \begin{longtable}{@{}>{\raggedright\arraybackslash}p{0.19\textwidth}>{\raggedright\arraybackslash}p{0.31\textwidth}>{\raggedright\arraybackslash}p{0.41\textwidth}@{}}%
  \toprule Source item & Modern mathematical content & Repository status \\\midrule\endhead%
}{\bottomrule\end{longtable}}
\title{The Haagerup--Zsid\'o \texorpdfstring{$\alpha=0$}{alpha=0} analytic bridge:\\
Poisson summation, odd-pole expansions, and the sharp reciprocal kernel}
\author{}
\date{Lean-oriented distilled synthesis for the finite Davis--Kahan formalization}
\begin{document}
\maketitle

\section*{Purpose and source map}
\begin{sourceguide}
\item[Primary analytic source]
Uffe Haagerup and L\'aszl\'o Zsid\'o, \emph{Resolvent Estimate for Hermitian Operators and a Related Minimal Extrapolation Problem}, Acta Sci. Math. (Szeged) 59 (1994), 503--524.
\item[Earlier $\alpha=0$ source]
B. Sz.-Nagy and A. Strausz, the 1938 extremal theorem cited as reference [14] by Haagerup--Zsid\'o, together with Sz.-Nagy's 1953 follow-up. The present note does not claim an independent reconstruction of those unavailable texts; it records the exact $\alpha=0$ theorem as stated and proved in the 1994 paper.
\item[Partial-fraction source]
The classical partial-fraction expansion of $\cothh$ and its use in accelerated Dirichlet-series formulas; a convenient modern reference is David M. Bradley, \emph{Series acceleration formulas for Dirichlet series with periodic coefficients}, Ramanujan J. 6 (2002), 331--346, arXiv:math/0505078.
\item[Formal Poisson source]
The pinned Mathlib modules \code{Mathlib.Analysis.Fourier.PoissonSummation}, \code{Mathlib.Analysis.Fourier.FourierTransform}, \code{Mathlib.Analysis.Fourier.Inversion}, and \code{Mathlib.Analysis.SpecialFunctions.ImproperIntegrals}.
\item[Repository target]
\code{ForMathlib/Analysis/Fourier/HaagerupZsidoKernel.lean} and the scalar root packaged by \code{HasIntegrableReciprocalFourierKernel}.
\item[Scope]
This is not a second source-order reconstruction of the 1994 paper. It is a focused synthesis of the smallest analytic theorem needed by the sharp finite-dimensional Sylvester and Davis--Kahan chain. Every displayed identity below is proved or reduced to explicitly named classical inputs.
\end{sourceguide}

The target is an integrable complex kernel $g:\R\to\C$ satisfying
\begin{equation}
  \int_{\R} g(t)e^{ixt}\,dt=\frac1x
  \qquad (\abs{x}\ge1),
  \tag{A0-Fourier}\label{A0-Fourier}
\end{equation}
and
\begin{equation}
  \int_{\R}\abs{g(t)}\,dt=\frac\pi2.
  \tag{A0-mass}\label{A0-mass}
\end{equation}
The 1994 paper proves a much larger minimal-extrapolation theorem. For the current formalization, \eqref{A0-Fourier} and \eqref{A0-mass} are the exact scalar root.

\section{The explicit $\alpha=0$ kernel}

\subsection{Fourier convention}

The source uses
\begin{equation}
  \widehat f(x)=\int_{\R}f(t)e^{ixt}\,dt.
  \tag{F}\label{F}
\end{equation}
This is an unnormalized Fourier transform with positive phase. Mathlib's notation $\mathcal F$ on $\R$ instead uses a $-2\pi ixt$ phase. The two conventions must never be conflated. The Mathlib transform is useful for Poisson summation; the final reciprocal identity must be translated back to \eqref{F}.

\subsection{Kernel formula}

Define the positive hyperbolic weight
\begin{equation}
  w(y):=\tanh\!\left(\frac{\pi y}{2}\right),
  \qquad y\ge0.
  \tag{weight}\label{weight}
\end{equation}
For $t\ne0$, define
\begin{equation}
  f_0(t)
  :=\frac{\sin t}{2}
     \int_0^\infty w(y)e^{-\abs{t}y}\,dy.
  \tag{kernel-real}\label{kernel-real}
\end{equation}
Assign any value at $t=0$; the natural Lean definition may use zero. The point value does not affect integrals.

Finally set
\begin{equation}
  g_0(t):=-i f_0(t).
  \tag{kernel-complex}\label{kernel-complex}
\end{equation}
The source proves
\begin{equation}
  \widehat f_0(x)=\frac{i}{x}
  \qquad(\abs{x}\ge1),
  \tag{f0-transform}\label{f0-transform}
\end{equation}
so \eqref{kernel-complex} gives \eqref{A0-Fourier}.

\begin{lemma}[Sign and parity]
The function $f_0$ is real-valued and odd. Moreover,
\[
  f_0(t)\sin t\ge0
  \qquad(t\ne0).
\]
Consequently
\[
  \abs{f_0(t)}
  =\frac{\abs{\sin t}}2
     \int_0^\infty w(y)e^{-\abs{t}y}\,dy.
\]
\end{lemma}

\begin{proofoutline}
\item The Laplace factor in \eqref{kernel-real} depends only on $\abs t$ and is nonnegative because $w(y)\ge0$ for $y\ge0$.
\item Multiplication by $\sin t$ therefore gives the sign assertion.
\item Since $\sin$ is odd and the Laplace factor is even, $f_0$ is odd.
\end{proofoutline}

\section{The exact mass: the shortest part of the proof}

The norm computation is independent of the exterior Fourier identity and should be formalized first.

\subsection{Laplace transform of $\abs{\sin}$}

\begin{lemma}[One-period integral]
For $y>0$,
\begin{equation}
  \int_0^\pi e^{-yt}\sin t\,dt
  =\frac{1+e^{-\pi y}}{1+y^2}.
  \tag{sin-period}\label{sin-period}
\end{equation}
\end{lemma}

\begin{proofoutline}
\item An antiderivative is
\[
  -\frac{e^{-yt}}{1+y^2}\bigl(y\sin t+\cos t\bigr).
\]
\item Evaluate at $0$ and $\pi$.
\end{proofoutline}

\begin{lemma}[Laplace transform of absolute sine]
For $y>0$,
\begin{equation}
  \int_0^\infty \abs{\sin t}e^{-yt}\,dt
  =\frac{1+e^{-\pi y}}{1-e^{-\pi y}}\frac1{1+y^2}
  =\frac{\cothh(\pi y/2)}{1+y^2}.
  \tag{abs-sin-laplace}\label{abs-sin-laplace}
\end{equation}
\end{lemma}

\begin{proofoutline}
\item Partition $\Ioi$ into $[n\pi,(n+1)\pi]$.
\item Write $t=n\pi+s$. Then $\abs{\sin(n\pi+s)}=\sin s$ for $0\le s\le\pi$.
\item The $n$th interval contributes
\[
  e^{-n\pi y}\int_0^\pi e^{-ys}\sin s\,ds.
\]
\item Sum the geometric series and use \eqref{sin-period}.
\end{proofoutline}

The formula is especially Lean-friendly because all summands are nonnegative. Tonelli and monotone convergence can be used without constructing an absolute-convergence argument for a signed series.

\subsection{Hyperbolic cancellation}

For $y>0$,
\begin{equation}
  \tanh\!\left(\frac{\pi y}{2}\right)
  \cothh\!\left(\frac{\pi y}{2}\right)=1.
  \tag{tanh-coth}\label{tanh-coth}
\end{equation}
If the available API for $\cothh$ is inconvenient, expand both functions using exponentials:
\[
  \tanh z=\frac{1-e^{-2z}}{1+e^{-2z}},
  \qquad
  \cothh z=\frac{1+e^{-2z}}{1-e^{-2z}}.
\]
For $z>0$, the denominators are nonzero and the product is visibly one.

\begin{theorem}[Exact $L^1$ mass]
The functions $f_0$ and $g_0$ are integrable and
\begin{equation}
  \int_{\R}\abs{f_0(t)}\,dt
  =\int_{\R}\abs{g_0(t)}\,dt
  =\frac\pi2.
  \tag{mass}\label{mass}
\end{equation}
\end{theorem}

\begin{proofoutline}
\item By the sign formula and Tonelli,
\begin{align*}
\int_{\R}\abs{f_0(t)}\,dt
&=\int_0^\infty
    w(y)\left(\int_0^\infty\abs{\sin t}e^{-yt}\,dt\right)dy.
\end{align*}
The factor $1/2$ in \eqref{kernel-real} cancels the factor two from the two halves of $\R$.
\item Insert \eqref{abs-sin-laplace} and cancel by \eqref{tanh-coth}:
\[
  \int_{\R}\abs{f_0(t)}\,dt
  =\int_0^\infty\frac{dy}{1+y^2}.
\]
\item Use the antiderivative $\arctan y$ and its limit $\pi/2$ at $+\infty$.
\item Since multiplication by $-i$ preserves the complex norm, $g_0$ has the same mass.
\end{proofoutline}

This proof simultaneously establishes integrability. A separate pointwise majorant is useful only if the formal Tonelli route becomes awkward.

\section{The elementary two-sided Cauchy transform}

For $y>0$ and $a\in\R$, define
\[
  C_y(a):=\int_{\R} e^{-y\abs t}e^{iat}\,dt.
\]

\begin{lemma}[Cauchy transform]
For $y>0$,
\begin{equation}
  C_y(a)=\frac{2y}{y^2+a^2}.
  \tag{cauchy-transform}\label{cauchy-transform}
\end{equation}
\end{lemma}

\begin{proofoutline}
\item Split $\R$ into $(-\infty,0]$ and $(0,\infty)$.
\item On the positive half-line integrate $e^{(-y+ia)t}$.
\item On the negative half-line substitute $u=-t$ and integrate $e^{(-y-ia)u}$.
\item Add
\[
  \frac1{y-ia}+\frac1{y+ia}=\frac{2y}{y^2+a^2}.
\]
\end{proofoutline}

\begin{corollary}[Inner sine transform]
For $y>0$ and $x\in\R$,
\begin{equation}
\int_{\R}\sin t\,e^{-y\abs t}e^{ixt}\,dt
=i y\left(
  \frac1{y^2+(x-1)^2}
  -\frac1{y^2+(x+1)^2}
\right).
\tag{sine-cauchy}\label{sine-cauchy}
\end{equation}
\end{corollary}

\begin{proofoutline}
\item Use
\[
  \sin t=\frac{e^{it}-e^{-it}}{2i}.
\]
\item Apply \eqref{cauchy-transform} at frequencies $x+1$ and $x-1$.
\item Check the sign carefully: the answer is $i$ times ``minus-frequency denominator first,'' as shown in \eqref{sine-cauchy}.
\end{proofoutline}

The agent's current status indicates that the Lean version of this checkpoint already compiles. It should remain the primitive transform lemma for both the Poisson and direct routes.

\section{Three routes to the exterior Fourier identity}

There are three mathematically sound routes. They have different formalization costs.

\subsection{Route A: source-faithful contour and limiting argument}

Haagerup--Zsid\'o obtain $f_0$ as the $L^1$ limit of their kernels $f_\alpha$ as $\alpha\to0$. For $\alpha\ne0$ they prove
\[
  \widehat f_\alpha(x)=\frac1{\alpha-ix}
  \qquad(\abs x\ge1).
\]
Continuity in $L^1$ implies uniform convergence of Fourier transforms, hence
\[
  \widehat f_0(x)=\frac1{-ix}=\frac{i}{x}.
\]

This route is source-faithful but expensive in Lean. It requires the full $\alpha$-family, contour deformation, dominated convergence in $L^1$, and endpoint control. It proves much more than the current project needs.

\subsection{Route B: Mittag--Leffler or product formula}

The zeros of $\cosh z$ occur at
\[
  z=i\pi\left(n+\frac12\right),\qquad n\in\mathbb Z.
\]
Taking the logarithmic derivative of the canonical product for $\cosh$ yields a partial-fraction expansion for $\tanh$. This is classical and concise on paper, but Mathlib does not currently supply the needed infinite product or Mittag--Leffler infrastructure in a directly usable form.

\subsection{Route C: Poisson summation and odd Cauchy poles}

This route uses APIs already present in Mathlib and is therefore the best current candidate. Its analytic heart is the Cauchy lattice sum, from which the odd-pole $\tanh$ expansion follows by subtracting the even lattice.

\section{Poisson summation and the Cauchy lattice}

\subsection{Normalized transform used only in this section}

For Poisson summation use Mathlib's normalized transform
\[
  (\FT h)(\xi)
  :=\int_{\R}e^{-2\pi i t\xi}h(t)\,dt.
\]
For $y>0$, set
\[
  h_y(t):=e^{-2\pi y\abs t}.
\]
By \eqref{cauchy-transform},
\begin{equation}
  (\FT h_y)(\xi)
  =\frac{y}{\pi(y^2+\xi^2)}.
  \tag{hy-transform}\label{hy-transform}
\end{equation}

\subsection{Poisson identity}

At $x=0$, Poisson summation gives
\[
  \sum_{n\in\mathbb Z}h_y(n)
  =\sum_{n\in\mathbb Z}(\FT h_y)(n).
\]
The left side is geometric:
\begin{equation}
  \sum_{n\in\mathbb Z}e^{-2\pi y\abs n}
  =1+2\sum_{n=1}^\infty e^{-2\pi yn}
  =\frac{1+e^{-2\pi y}}{1-e^{-2\pi y}}
  =\cothh(\pi y).
  \tag{geom-coth}\label{geom-coth}
\end{equation}
Using \eqref{hy-transform} yields the classical Cauchy lattice formula.

\begin{theorem}[Cauchy lattice sum]
For $y>0$,
\begin{equation}
  \sum_{n\in\mathbb Z}\frac1{y^2+n^2}
  =\frac\pi y\cothh(\pi y).
  \tag{cauchy-lattice}\label{cauchy-lattice}
\end{equation}
\end{theorem}

\subsection{Extracting the odd lattice}

The even part is
\begin{align}
  \sum_{n\in\mathbb Z}\frac1{y^2+(2n)^2}
  &=\frac14
    \sum_{n\in\mathbb Z}\frac1{(y/2)^2+n^2}
  \notag\\
  &=\frac\pi{2y}\cothh\!\left(\frac{\pi y}{2}\right).
  \tag{even-lattice}\label{even-lattice}
\end{align}
Subtract \eqref{even-lattice} from \eqref{cauchy-lattice}. The elementary identity
\begin{equation}
  \cothh(2z)-\frac12\cothh z
  =\frac12\tanh z
  \tag{double-coth}\label{double-coth}
\end{equation}
gives
\begin{equation}
  \sum_{\substack{k\in\mathbb Z\\ k\text{ odd}}}
    \frac1{y^2+k^2}
  =\frac\pi{2y}\tanh\!\left(\frac{\pi y}{2}\right).
  \tag{odd-lattice}\label{odd-lattice}
\end{equation}
The summand is even in $k$, so the positive odd terms contribute half.

\begin{theorem}[Odd-pole expansion of $\tanh$]
For $y>0$,
\begin{equation}
  \frac{\tanh(\pi y/2)}{y}
  =\frac4\pi\sum_{n=0}^\infty
    \frac1{y^2+(2n+1)^2}.
  \tag{tanh-odd}\label{tanh-odd}
\end{equation}
\end{theorem}

This is the exact scalar series needed below. It is also the piece of classical analysis most worth isolating as a reusable `ForMathlib` theorem.

\section{The rational integral that makes the series telescope}

\begin{lemma}[Two Cauchy denominators]
Let $a\ge0$ and $c>0$. Then
\begin{equation}
  \int_0^\infty
  \frac{y^2}{(y^2+c^2)(y^2+a^2)}\,dy
  =\frac\pi{2(a+c)}.
  \tag{two-cauchy}\label{two-cauchy}
\end{equation}
\end{lemma}

\begin{proofoutline}
\item If $a=0$, the integrand is $(y^2+c^2)^{-1}$ and the result is immediate.
\item If $a>0$ and $a\ne c$, use
\[
\frac{y^2}{(y^2+c^2)(y^2+a^2)}
=\frac{c^2}{c^2-a^2}\frac1{y^2+c^2}
 -\frac{a^2}{c^2-a^2}\frac1{y^2+a^2}.
\]
\item Integrate each Cauchy kernel:
\[
  \int_0^\infty\frac{dy}{y^2+r^2}=\frac\pi{2r}
  \qquad(r>0).
\]
\item If $a=c$, either take the limit or prove directly using the derivative of $y/(y^2+c^2)$, or use
\[
  \frac{y^2}{(y^2+c^2)^2}
  =\frac1{y^2+c^2}-\frac{c^2}{(y^2+c^2)^2}.
\]
\end{proofoutline}

For the application, only $c=2n+1$ and $a=x-1\ge0$ occur. Thus $c>0$ is automatic and the delicate $a=c$ branch occurs at isolated real values of $x$; nevertheless the clean general lemma should handle it.

\section{The telescoping identity}

\begin{theorem}[Scalar exterior multiplier identity]
For $a\ge0$,
\begin{equation}
\begin{aligned}
&\int_0^\infty
  \tanh\!\left(\frac{\pi y}{2}\right)y
  \left(
    \frac1{y^2+a^2}
    -\frac1{y^2+(a+2)^2}
  \right)dy\\
&\hspace{5cm}=\frac2{a+1}.
\end{aligned}
\tag{telescoping-integral}\label{telescoping-integral}
\end{equation}
\end{theorem}

\begin{proofoutline}
\item Insert \eqref{tanh-odd}:
\[
  \tanh\!\left(\frac{\pi y}{2}\right)y
  =\frac4\pi\sum_{n=0}^\infty
     \frac{y^2}{y^2+(2n+1)^2}.
\]
\item For $a\ge0$, the bracket in \eqref{telescoping-integral} is nonnegative. Every series term is therefore nonnegative. Tonelli or monotone convergence justifies exchanging sum and integral.
\item Apply \eqref{two-cauchy} twice with $c_n=2n+1$:
\begin{align*}
&\int_0^\infty
\frac{y^2}{y^2+c_n^2}
\left(\frac1{y^2+a^2}-\frac1{y^2+(a+2)^2}\right)dy\\
&\qquad=\frac\pi2
\left(\frac1{a+c_n}-\frac1{a+2+c_n}\right).
\end{align*}
\item Multiplication by $4/\pi$ gives
\[
  2\sum_{n=0}^\infty
  \left(
    \frac1{a+2n+1}-\frac1{a+2n+3}
  \right).
\]
\item The $N$th partial sum is
\[
  2\left(\frac1{a+1}-\frac1{a+2N+3}\right),
\]
which converges to $2/(a+1)$.
\end{proofoutline}

This telescoping is the main conceptual payoff of the Poisson route. The infinite harmonic-analysis expression collapses to one boundary term.

\section{Exterior Fourier identity}

\begin{theorem}[Positive exterior frequencies]
For $x\ge1$,
\begin{equation}
  \widehat f_0(x)=\frac{i}{x}.
  \tag{positive-exterior}\label{positive-exterior}
\end{equation}
\end{theorem}

\begin{proofoutline}
\item Substitute \eqref{kernel-real} into the Fourier integral.
\item Use Fubini. The absolute double integral is bounded by the mass calculation, so no new domination theorem is needed.
\item Apply \eqref{sine-cauchy}:
\begin{align*}
\widehat f_0(x)
&=\frac i2\int_0^\infty w(y)y
\left(
  \frac1{y^2+(x-1)^2}
 -\frac1{y^2+(x+1)^2}
\right)dy.
\end{align*}
\item Set $a=x-1\ge0$. Then $x+1=a+2$.
\item Apply \eqref{telescoping-integral} to obtain
\[
  \widehat f_0(x)=\frac i2\frac2{a+1}=\frac i x.
\]
\end{proofoutline}

\begin{theorem}[Negative exterior frequencies]
For $x\le-1$,
\begin{equation}
  \widehat f_0(x)=\frac{i}{x}.
  \tag{negative-exterior}\label{negative-exterior}
\end{equation}
\end{theorem}

\begin{proofoutline}
\item Since $f_0$ is real and odd, its Fourier transform under convention \eqref{F} is imaginary and odd.
\item Apply \eqref{positive-exterior} to $-x\ge1$.
\end{proofoutline}

\begin{corollary}[Sharp reciprocal kernel]
The complex kernel $g_0=-if_0$ satisfies
\[
  \widehat g_0(x)=\frac1x\qquad(\abs x\ge1)
\]
and
\[
  \norm{g_0}_{L^1}=\frac\pi2.
\]
\end{corollary}

\section{Lean-sized theorem seams}

The following decomposition minimizes simultaneous elaboration risk.

\begin{formalizationmap}
Hyperbolic weight & $0\le w(y)\le1$ for $y\ge0$; exponential formula for $w$; $w\cdot\cothh=1$ for $y>0$. & New scalar helpers in \code{HaagerupZsidoKernel.lean}.\\
Two-sided Laplace transform & Equation \eqref{cauchy-transform} and sine corollary \eqref{sine-cauchy}. & Agent reports this checkpoint already compiles.\\
Absolute sine Laplace transform & Equations \eqref{sin-period} and \eqref{abs-sin-laplace}. & Prefer nonnegative interval decomposition plus geometric series.\\
Exact mass & Equation \eqref{mass}. & Use Tonelli; this proves integrability and norm equality together.\\
Normalized exponential transform & Equation \eqref{hy-transform}. & Reuse two-sided Laplace transform with the $2\pi$ normalization made explicit.\\
Cauchy lattice & Equation \eqref{cauchy-lattice}. & Apply Mathlib Poisson summation to $h_y(t)=e^{-2\pi y\abs t}$.\\
Odd-pole $\tanh$ expansion & Equation \eqref{tanh-odd}. & Split the integer lattice into even and odd terms; avoid residues.\\
Rational integral & Equation \eqref{two-cauchy}. & Isolate zero, equal-parameter, and generic partial-fraction branches.\\
Telescoping integral & Equation \eqref{telescoping-integral}. & All terms nonnegative; exchange sum and integral by Tonelli; prove finite telescoping explicitly.\\
Exterior transform & Equations \eqref{positive-exterior} and \eqref{negative-exterior}. & Positive branch from telescoping; negative branch from real odd symmetry.\\
Kernel package & Equations \eqref{A0-Fourier} and \eqref{A0-mass}. & Construct the \texttt{\seqsplit{HasIntegrableReciprocalFourierKernel}} certificate at mass $\pi/2$.\\
\end{formalizationmap}

Suggested declaration shapes:

\begin{verbatim}
theorem integral_abs_sin_mul_exp_neg
    {y : R} (hy : 0 < y) :
    integral (Ioi 0) (fun t => |sin t| * exp (-y * t))
      = ((1 + exp (-pi * y)) / (1 - exp (-pi * y)))
          / (1 + y^2)

 theorem cauchy_lattice_sum
    {y : R} (hy : 0 < y) :
    tsum (fun n : Z => (y^2 + (n : R)^2)^(-1))
      = pi / y * coth (pi * y)

 theorem tanh_div_eq_odd_cauchy_tsum
    {y : R} (hy : 0 < y) :
    tanh (pi * y / 2) / y
      = (4 / pi) * tsum (fun n : N =>
          (y^2 + ((2*n+1 : N) : R)^2)^(-1))

 theorem integral_tanh_reciprocal_difference
    {a : R} (ha : 0 <= a) :
    integral (Ioi 0) (fun y =>
      tanh (pi*y/2) * y *
        ((y^2+a^2)^(-1) - (y^2+(a+2)^2)^(-1)))
      = 2 / (a+1)
\end{verbatim}

These are schematic, not guaranteed exact syntax. Preserve the mathematical boundaries even if the pinned APIs suggest different names.

\section{Failure modes and route-selection advice}

\begin{enumerate}[leftmargin=2em]
\item \textbf{Normalization drift.} Poisson summation uses a $2\pi$-normalized transform; the final reciprocal theorem does not. State conversion lemmas rather than relying on simplification.
\item \textbf{Continuity at zero.} The kernel need not be continuous at zero. Do not make continuity a prerequisite for integrability or Fourier evaluation.
\item \textbf{Signed Fubini too early.} Prove the exact nonnegative mass identity first. It supplies the absolute-integrability certificate needed by the transform calculation.
\item \textbf{Integer parity bookkeeping.} Prove a dedicated equivalence between odd integers and $\pm(2n+1)$. Avoid repeatedly normalizing casts inside a large calculation.
\item \textbf{Infinite telescoping by automation.} Prove the finite partial-sum identity and then take a limit. `simp` should not be asked to discover the telescoping structure.
\item \textbf{Exact finite orbit attainment.} The $L^1$ kernel theorem is enough for sharp inequalities through $\pi/2+\varepsilon$ finite approximants. Do not block the analytic theorem on a stronger compactness statement.
\item \textbf{Overformalizing minimal extrapolation.} Weak-star compactness, uniqueness of minimizers, and the full $\alpha$-family are valuable future contributions but not prerequisites for the finite DK theorem.
\end{enumerate}

\section{What the literature contributes to the current proof}

Haagerup--Zsid\'o supply the explicit extremal kernel, its sign structure, exact mass, and exterior transform. Their theorem confirms that the desired $\pi/2$ constant is not a convenient upper bound but the optimal minimal extrapolation norm.

The classical $\cothh$ partial fraction identity supplies an alternative derivation of the only difficult scalar transform. Bradley's exposition is useful because it treats the partial fraction expansion as a reusable analytic engine rather than an isolated table entry.

Mathlib's Poisson summation theorem makes that classical identity reachable without first developing general Mittag--Leffler theory. The exponential Cauchy kernel has elementary decay and an elementary Fourier transform, exactly matching the hypotheses for the formal Poisson API.

Thus the literature synthesis supports a proof strategy that is simultaneously:
\begin{itemize}
\item faithful to the sharp Haagerup--Zsid\'o theorem;
\item substantially smaller than their full contour/minimal-extrapolation proof;
\item based on analysis already represented in Mathlib;
\item modular enough to preserve useful scalar results even if the final exchange of sum and integral needs iteration.
\end{itemize}

\section*{Recommended implementation order}

\begin{enumerate}[leftmargin=2em]
\item Keep the compiled hyperbolic-weight and Cauchy-transform checkpoint.
\item Prove \eqref{abs-sin-laplace} and the exact mass theorem.
\item Prove the normalized Fourier transform \eqref{hy-transform}.
\item Apply Poisson summation to obtain \eqref{cauchy-lattice}.
\item Derive the odd-pole expansion \eqref{tanh-odd}.
\item Prove the rational integral \eqref{two-cauchy}.
\item Prove the telescoping identity \eqref{telescoping-integral}.
\item Complete the positive exterior transform, then obtain the negative branch by oddness.
\item Package the reciprocal kernel and immediately run the axiom audit on the real and complex Ky Fan and Sylvester endpoints.
\end{enumerate}

\section*{Bibliographic references}

\begin{thebibliography}{9}
\bibitem{HZ1994}
U. Haagerup and L. Zsid\'o,
\emph{Resolvent estimate for Hermitian operators and a related minimal extrapolation problem},
Acta Sci. Math. (Szeged) 59 (1994), 503--524.

\bibitem{Bradley2002}
D. M. Bradley,
\emph{Series acceleration formulas for Dirichlet series with periodic coefficients},
Ramanujan J. 6 (2002), 331--346; arXiv:math/0505078.

\bibitem{SzNagyStrausz1938}
B. Sz.-Nagy and A. Strausz,
\emph{On a theorem of H. Bohr} (Hungarian),
Mat. Term\'eszettud. \`Ertes\'it\H{o} 37 (1938), 121--133.

\bibitem{SzNagy1953}
B. Sz.-Nagy,
\emph{\"Uber die Ungleichung von H. Bohr},
Math. Nachr. 9 (1953), 255--259.
\end{thebibliography}

\section*{Source-faithfulness statement}

The explicit kernel and the sharp mass and exterior-transform claims are reconstructed from Haagerup--Zsid\'o's $\alpha=0$ theorem. The Poisson-summation derivation of the odd-pole $\tanh$ expansion is a modern alternative proof assembled for formalization; it is not presented as the proof used in the 1994 article. The note deliberately distinguishes source content from the proposed Lean route.

\end{document}
